In distribution centres, the order picking problem ({\bf OPP}) consists of assigning batches of orders to pickers, determining the order
in which each picker executes them, and establishing the picking route through the warehouse. The objective is to minimise the sum of the
batch completion times, together with the penalties for earliness and tardiness with respect to the committed due dates. Although solving this
problem exactly guarantees optimality, the computation time grows quickly with the size of the warehouse and the number of orders. This
hinders its use in real operations, where decisions must be made within limited response times.

This work takes as its starting point a mixed integer linear programming model for the OPP from the literature. Using it as a guide, we
formulate our own linear model, which we call {\em monolithic}, that incorporates heterogeneous picker speeds and cold chain constraints.
Then, in order to reduce its running times and address larger instances, this model is decomposed into two stages that are solved one after
the other: first the routing of each batch, and then the assignment and sequencing of the batches. We call this second model {\em Split}.
Since the Split model is still too costly to solve in practical times for moderately large instances, two metaheuristics based on
{\em Hill Climbing} and {\em Ant Colony Optimization} are developed, equivalent to the linear models in the features they consider, except for voluntary waiting.

The four methods are evaluated on 38 functional test cases and 19 stress scenarios, most of them with a realistic warehouse profile. The
Split model reaches the same optimum as the monolithic model and is the only one of the two that proves it at a realistic scale, in 9 of
the 19 stress scenarios. In the main runs, the metaheuristics find a valid solution in every
feasible instance. In the stress tests, {\em Ant Colony Optimization}
stays within 0.04\,\% of the optimum and {\em Hill Climbing} within 0.21\,\%. {\em Hill Climbing} responds about 53 times faster
than the Split in the median, while {\em Ant Colony Optimization} takes from slightly longer than the Split to 6 times less,
depending on the scenario. Since no method dominates the others in quality, time and robustness at once, we conclude that the Split is preferable when
a proven optimum is needed, in instances of the realistic profile with up to 6 batches on grids of up to 15×15; {\em Ant Colony Optimization}, when solution quality
matters, and {\em Hill Climbing}, when a response within seconds is required.
